% ======================================================================
% SECTION 2: AUTONOMOUS SPONSOR GOVERNANCE AND PORTFOLIO MANAGEMENT
% File: sections/governance.tex
% ======================================================================

The governance layer replaces the traditional board of directors, portfolio committee, and
asset-level leadership with autonomous agents that execute strategic decisions through
algorithmic logic. Under 21 CFR \S312.50, the sponsor is responsible for selecting qualified
investigators, providing required information, monitoring the investigation, and ensuring
compliance with regulatory requirements \cite{cfr-part-312}. ICH E6(R3) further requires a
quality-by-design approach to clinical trial governance, with prospective identification of
critical-to-quality factors and systematic risk management \cite{ich-e6r3}. The governance
layer implements these requirements through three agent subsystems that collectively manage
portfolio strategy, asset-level decision-making, and lifecycle gate progression.

\subsection{Board Agent Cluster}

The board agent cluster replaces the executive committee and study steering committee for
trial governance decisions. It operates through three interconnected modules: charter
management, risk appetite quantification, and capital allocation optimization.

The charter management module maintains the operating charter for each active clinical
program. This charter defines the therapeutic area focus, target patient populations,
competitive positioning requirements, and success criteria at each development stage.
When new clinical or competitive intelligence arrives, the charter management module
evaluates whether the existing charter remains valid or requires revision. Charter revisions
above a defined materiality threshold (changes affecting primary endpoints, target populations,
or investment commitments exceeding 15\% of program budget) trigger a formal review cycle
with documented rationale.

The risk appetite quantification module translates qualitative strategic intent into
quantitative decision boundaries. For each program, it defines acceptable ranges for
probability of technical success, probability of regulatory success, expected net present
value, and time-to-market. These boundaries create an envelope within which the portfolio
agent and asset lead agent operate autonomously. When any program metric moves outside the
defined envelope, the board agent cluster either adjusts boundaries (within its own authority
for minor deviations) or escalates to the human sponsor-of-record for review.

The capital allocation optimizer distributes investment across the portfolio to maximize
risk-adjusted returns. It considers phase transition probabilities, competitive timing
pressures, and regulatory pathway advantages. The optimizer runs continuously as new data
arrives from ongoing trials, generating rebalancing recommendations when portfolio
composition deviates from target allocations by more than a configurable threshold.

When algorithmic confidence drops below 85\% on any major decision (defined as decisions
affecting patient safety, primary endpoint changes, or investment commitments exceeding
\$5 million), the board agent cluster escalates to the human sponsor-of-record. This
threshold is calibrated to ensure autonomous operation for routine portfolio management
while preserving human oversight for high-stakes decisions. The escalation policy
generates a structured briefing document containing the decision context, agent
recommendation, confidence score, risk analysis, and recommended alternatives.

\subsection{Portfolio Governance Agent}

The portfolio governance agent (portfolio\_agent) replaces the research and development
portfolio governance function. It automates three core activities: program prioritization
using clinical differentiation scoring, indication strategy selection based on competitive
landscape analysis, and milestone-triggered decision-making across the development lifecycle.

Program prioritization uses a composite scoring model that integrates clinical differentiation
potential, biomarker prevalence and testing access feasibility, dose optimization readiness
(per FDA Project Optimus requirements \cite{fda-project-optimus}), endpoint strategy
sustainability across approval pathways, and commercial viability. Each scoring dimension is
weighted according to the portfolio strategy and therapeutic area priorities. Programs are
ranked by composite score, and resource allocation follows the ranking subject to portfolio
balance constraints (such as maintaining a target distribution across early-stage, mid-stage,
and late-stage assets).

Indication strategy selection evaluates the competitive landscape for each molecular target
and therapeutic modality. The agent tracks the evolving modality mix, including antibody-drug
conjugates, multispecific antibodies, and cell and gene therapies, and assesses how emerging
modalities affect portfolio positioning. For each indication under consideration, the agent
generates a competitive assessment covering mechanism differentiation, patient population
overlap with existing therapies, regulatory pathway availability (including accelerated
approval potential), and projected time-to-market relative to competitors.

The milestone-triggered decision framework maps to the seven decision gates defined for
oncology drug development \cite{tufts-csdd-cost}. At each gate, the portfolio agent evaluates
whether the program meets predefined go-criteria (target validation strength, dossier
completeness, dose-exposure-response characterization, endpoint agreement, enrollment
tracking, data quality within tolerance, and confirmatory commitment readiness). Programs
that meet all go-criteria advance automatically. Programs that partially meet criteria enter
a conditional-advance state with specific remediation requirements. Programs that fail
critical criteria trigger a hold recommendation with documented rationale.

The biomarker funnel math module calculates the operational feasibility of biomarker-driven
trials. Given a target biomarker, it computes the effective patient population by
multiplying patient availability by testing coverage, biomarker prevalence, trial
eligibility rate, and consent rate. This calculation identifies programs where biomarker
selection creates enrollment bottlenecks, enabling portfolio decisions that account for
operational feasibility alongside scientific merit.

\subsection{Asset Lead Agent}

The asset lead agent (asset\_lead\_agent) replaces the clinical development head or asset lead
for individual programs. It operates at the intersection of the portfolio layer and the study
execution layer, translating portfolio strategy into study-level execution plans.

The agent generates asset strategy memos that document the target product profile, competitive
landscape positioning, endpoint strategy, biomarker development plan, dose justification
(including Project Optimus compliance \cite{fda-project-optimus}), and lifecycle plan. These
memos serve as the governing document for each program and are updated automatically when
new clinical data, competitive intelligence, or regulatory guidance becomes available.

Vendor strategy automation selects and manages the contract research organization and
specialty vendor relationships for each study. The agent evaluates CRO capabilities against
study requirements (therapeutic area expertise, geographic coverage, technology platform
compatibility, and historical performance metrics), generates requests for proposal,
evaluates responses against predefined criteria, and manages ongoing vendor relationships
through service-level agreement monitoring. The vendor strategy module interfaces with the
vendor management agent mesh (\S\ref{sec:vendor}) for operational vendor oversight.

Study-level KPI accountability tracks execution performance against the decision gate
criteria. The asset lead agent monitors enrollment velocity, screen failure rates, protocol
deviation rates, data query resolution times, and safety reporting timeliness. When KPIs
deviate from targets, the agent generates corrective action recommendations and, if
deviations persist, escalates to the portfolio agent for portfolio-level assessment.

\subsection{Governance Decision Framework}

The governance decision framework provides the structural foundation for autonomous
decision-making across all twelve agents. Table~\ref{tab:agent-mapping} maps each agent to
its traditional human role equivalent, key functions, and decision authority scope.

\begin{longtable}{L{2.2cm} L{2.0cm} L{2.0cm} L{2.8cm} L{2.8cm}}
\caption{Governance Agent Mapping: Traditional Roles to Autonomous Agents}
\label{tab:agent-mapping} \\
\toprule
\textbf{Agent Name} & \textbf{Traditional Role} & \textbf{Layer} & \textbf{Key Functions} & \textbf{Decision Authority} \\
\midrule
\endfirsthead
\toprule
\textbf{Agent Name} & \textbf{Traditional Role} & \textbf{Layer} & \textbf{Key Functions} & \textbf{Decision Authority} \\
\midrule
\endhead
\midrule
\multicolumn{5}{r}{\textit{Continued on next page}} \\
\endfoot
\bottomrule
\endlastfoot
portfolio\_agent & VP R\&D, Portfolio Committee & Governance & Program prioritization, indication strategy, milestone gates & Go/no-go up to threshold; escalate above \\
asset\_lead\_agent & Clinical Development Head & Governance & Asset strategy, vendor selection, KPI tracking & Study-level decisions within portfolio bounds \\
clinical\_accountability\_agent & Study Physician, Clinical Lead & Execution & Protocol logic, eligibility, dose optimization & Medical decisions with safety escalation \\
study\_orchestrator & Cross-functional Study Team & Execution & Dependency management, work routing, issue escalation & Coordination decisions, no clinical authority \\
clinops\_agent & Clinical Operations Director & Execution & Site startup, monitoring, enrollment, query resolution & Operational decisions within protocol \\
safety\_agent & Safety and Pharmacovigilance Team & Execution & AE intake, signal detection, expedited reporting & Safety classification with human gate for SAEs \\
regulatory\_agent & Regulatory Affairs Director & Execution & IND lifecycle, filing, authority correspondence & Filing decisions with sponsor-of-record review \\
quality\_agent & Quality Assurance Director & Execution & SOP enforcement, CAPA, audit readiness, RBQM & Quality decisions with escalation for critical findings \\
supply\_agent & CMC and Supply Chain Lead & Execution & IMP logistics, demand forecasting, lot release & Supply decisions within validated parameters \\
data\_biostats\_agent & Data Management and Biostatistics & Execution & Data pipeline, CDISC, SAP, interim analysis & Data decisions; unblinding requires human gate \\
site\_gateway & Site Relationship Manager & Site/Robotics & Site qualification, training, scheduling & Site-level operational decisions \\
robot\_execution\_gateway & (No traditional equivalent) & Site/Robotics & Task orders, safety gates, procedure state & Procedure authorization with mandatory human gates \\
\end{longtable}

Table~\ref{tab:decision-gates} presents the seven decision gates adapted from the
oncology drug development lifecycle \cite{tufts-csdd-cost}, showing the go-criteria,
KPI thresholds, agent ownership, and escalation rules for each gate.

\begin{longtable}{L{2.0cm} L{2.8cm} L{2.5cm} L{2.0cm} L{2.5cm}}
\caption{Decision Gate Framework: Seven Lifecycle Gates with Autonomous Trigger Conditions}
\label{tab:decision-gates} \\
\toprule
\textbf{Gate} & \textbf{Go-Criteria} & \textbf{KPI Thresholds} & \textbf{Agent Owner} & \textbf{Escalation Rule} \\
\midrule
\endfirsthead
\toprule
\textbf{Gate} & \textbf{Go-Criteria} & \textbf{KPI Thresholds} & \textbf{Agent Owner} & \textbf{Escalation Rule} \\
\midrule
\endhead
\midrule
\multicolumn{5}{r}{\textit{Continued on next page}} \\
\endfoot
\bottomrule
\endlastfoot
Candidate Selection & Target validation, translational biomarker plan, manufacturability & Probability of technical success $\geq$ 40\% & portfolio\_agent & Escalate if confidence $<$ 85\% \\
IND/CTA Green Light & Dossier complete, assay plan locked, safety monitoring plan & Dossier completeness $\geq$ 95\%, assay validated & regulatory\_agent & Escalate if any critical component missing \\
Dose Optimization Readiness & Dose-exposure-response characterized, tolerability confirmed & Therapeutic index within target range & clinical\_accountability\_agent & Mandatory human review per Project Optimus \\
Pivotal Design Lock & Endpoint and estimand agreed, SAP validated, CDx path feasible & Regulatory alignment confirmed & clinical\_accountability\_agent & Escalate if endpoint not agreed with authority \\
Pivotal Execution Control & Enrollment on track, data quality within tolerance, safety manageable & Enrollment $\geq$ 80\% of target velocity, QTLs within limits & study\_orchestrator & Escalate if $\geq$ 2 KRIs breached simultaneously \\
Database Lock and CSR & Data cleaned, audit trail complete, CSR aligned with ICH E3 & Query resolution $\geq$ 99\%, no open critical queries & data\_biostats\_agent & Escalate if audit trail gaps detected \\
Launch and Post-marketing & Confirmatory commitments funded, access dossier ready & Submission package validated, disclosure timeline confirmed & regulatory\_agent & Mandatory sponsor-of-record sign-off \\
\end{longtable}

The escalation mechanism ensures that autonomous operation is bounded by human oversight at
critical junctures. Three categories of escalation exist. Confidence-based escalation occurs
when any agent's algorithmic confidence score drops below 85\% on a decision within its
scope. Mandatory-gate escalation applies to predefined decision types that always require
human review: initial IND filing, primary endpoint changes, safety signal escalation to
regulatory authorities, database lock authorization, and robotic procedure authorization for
Tier 3 (highest risk) procedures. Override escalation occurs when the human sponsor-of-record
proactively intervenes to review or reverse any agent decision, which is an unrestricted right
under the system's operating charter.
